% !TeX program = xelatex
% UTF-8 transcription, in the order of the five original photographs.
% pic01.jpg = 20261012_225122644.jpg
% pic02.jpg = 20261012_225136733.jpg
% pic03.jpg = 20261012_225143258.jpg
% pic04.jpg = 20261012_225155655.jpg
% pic05.jpg = 20261012_225203159.jpg
% Medical statements and abbreviations are retained as written in the source.
% Compile directly with XeLaTeX, or input into a XePersian book.
\newif\ifMSStandalone
\ifdefined\settextfont
  \MSStandalonefalse
\else
  \MSStandalonetrue
\fi
\ifMSStandalone
  \documentclass[a4paper,12pt]{extbook}
  \usepackage[margin=2.5cm]{geometry}
  \usepackage{xcolor}
  \usepackage{amsmath,amssymb}
  \usepackage{enumitem}
  \usepackage{xepersian}
  \settextfont{Amiri}
  \setlatintextfont{FreeSerif}
  \pagestyle{plain}
\fi
\ifMSStandalone
  \begin{document}
\fi
\begingroup
\emergencystretch=2em
\newcommand{\MSen}[1]{\lr{#1}}
\ifdefined\colorbox
  \newcommand{\MShl}[1]{\colorbox{yellow!35}{\strut #1}}
\else
  \newcommand{\MShl}[1]{\textbf{#1}}
\fi
\newcommand{\MSpage}[2]{%
  \clearpage\section*{صفحهٔ #1 ـ \MSen{MS} ـ ۱۳}%
  \noindent{\footnotesize\MSen{#2}}\par\smallskip
  {\footnotesize\MSen{Date: / /}\quad\MSen{Subject: MS -- 13}\par
  \MSen{Sat.\quad Sun.\quad Mon.\quad Tue.\quad Thu.\quad Wed.\quad Fri.}}\par\medskip}

\chapter{\MSen{MS}}
\MSpage{۱}{pic01.jpg}
\begin{itemize}
\item \MSen{MS}: میلین توسط الیگودندروسیت‌ها تولید می‌شود.
\item \MShl{در زنان ۲ برابر} شایع‌تر / سن شایع \MShl{۲۰–۴۰ سالگی}.
\item کمترین شیوع: مناطق استوایی / بیشترین: در شمال اسکاتلند.
\item مواجهه با عامل محیطی قبل از ۱۵ سالگی ریسک ابتلا را بالا می‌برد.
\item در سفیدپوستان شایع‌تر / در خواهر و برادر بیماران ۵–۱۵ برابر شایع‌تر.
\item در دوقلوهای مونوزیگوت بیشتر.
\item \MSen{MHC} روی کروموزوم ۶ / \MShl{\MSen{HLA-DR2}} / \MShl{\MSen{HLA-DW2}}.
\item برخی ویروس‌ها: سرخک، تبخال، زونا، آنفولانزا \MSen{C}، \MSen{EBV}، \MSen{HPV6}، کلامیدیا پنومونیه.
\item هم لنفوسیت \MSen{B} هم \MSen{T}: \MShl{واکنش سلولی و هومورال}.
\item درون پلاک‌ها، سلول‌های لنفوسیت \MSen{B} و \MSen{T} و ماکروفاژ.
\item انهدام الیگودندروسیت‌ها در پلاک‌های مزمن.
\item گلیوز و پرولیفراسیون آستروسیت‌ها.
\end{itemize}
\section*{علائم}
\begin{itemize}
\item \MShl{نوریت اپتیک}: کاهش بینایی، تاری دید، اختلال دید رنگ قرمز، درد در حرکات.
\begin{itemize}
\item اکثراً یک‌طرفه / موقتی / یک تا چند هفته / دیسک اپتیک در آغاز \MSen{NL}.
\item اگر مزمن: آتروفی عصب بینایی.
\end{itemize}
\item \MShl{ضعف اندام‌ها}: اندام‌های تحتانی یا هر چهار اندام یا یک اندام.
\begin{itemize}
\item \MShl{ضعف} با سایر \MShl{علائم درگیری پیرامیدال} همراه است: \MShl{اسپاستیته}، \MShl{\MSen{$DTR\uparrow$}}، علامت \MShl{بابنسکی}.
\end{itemize}
\item \MShl{خستگی زودرس}: در بیش از ۹۰٪ بیماران.
\item علائم حسی: پارستزی / احساس تورم و سنگینی / درد / کاهش یا از بین رفتن حس سطحی یا عمقی.
\item اگر سطح حسی وجود داشته باشد: ضایعه در نخاع.
\end{itemize}

\MSpage{۲}{pic02.jpg}
\begin{itemize}
\item \MShl{دوبینی}: اکثراً به علت افتالموپلژی بین‌هسته‌ای (\MShl{\MSen{INO}}) یا \MShl{فلج زوج ۶}.
\begin{itemize}
\item \MSen{INO} به علت آسیب به \MShl{\MSen{MLF}} ایجاد می‌شود.
\item در \MSen{INO}، \MShl{چشم سمت درگیر نمی‌تواند به داخل نگاه کند} و چشم مقابل در نگاه به خارج دچار نیستاگموس می‌شود.
\item شایع‌ترین علت \MSen{INO} در افراد جوان: \MShl{\MSen{MS}}.
\end{itemize}
\item آتاکسی: به علت ایجاد پلاک در \MShl{مخچه}، \MShl{\MSen{brainstem}} یا \MShl{نخاع}.
\item \MShl{اختلال اسفنکتری}: در بیش از ۹۰٪ بیماران، به صورت ارجانسی، بی‌اختیاری، تکرر.
\begin{itemize}
\item \MShl{یبوست}.
\end{itemize}
\item مشکلات شناختی: حافظه، تمرکز، تجزیه و تحلیل.
\item \MShl{افسردگی}: در بیش از ۵۰٪ بیماران / \MShl{خودکشی} در \MSen{MS} شایع است.
\item عدم تحمل گرما: با افزایش دما علائم بیماری تشدید می‌شود.
\begin{itemize}
\item \MShl{پدیدهٔ \MSen{Uhthoff}}.
\end{itemize}
\item \MShl{علامت لرمیت}: \MShl{نشانهٔ دمیلینهٔ نخاع} / با خم کردن گردن احساس برق‌گرفتگی در [خوانش نامطمئن: تنه] یا اندام‌ها.
\item سرگیجه، اختلالات جنسی، نورالژی تری‌ژمینال، فلج عصب ۷.
\end{itemize}
\section*{انواع سیر بالینی \MSen{MS}}
\begin{itemize}
\item \MShl{عودکننده ـ بهبود یابنده}: ۸۵٪ موارد / در بین حملات بیماری پیشرفت نمی‌کند.
\begin{itemize}
\item حمله باید \MShl{حداقل \MSen{24 h}} ادامه یابد و \MShl{بین حملات حداقل ۱ ماه فاصله} باشد.
\end{itemize}
\item \MShl{پیش‌روندهٔ ثانویه}: در طول زمان با یا بدون وجود حمله پیشرفت می‌کند.
\begin{itemize}
\item \MShl{بیشتر موارد \MSen{MS} عودکننده} به \MShl{این نوع تبدیل می‌شوند}.
\end{itemize}
\item پیش‌روندهٔ اولیه: بدون وجود حمله از ابتدا پیشرفت می‌کند / به سرعت دچار ناتوانی.
\item پیش‌رونده ـ عودکننده: در فاصلهٔ بین حملات هم پیشرفت می‌کند.
\end{itemize}

\MSpage{۳}{pic03.jpg}
\section*{تشخیص \MSen{MS}}
\begin{itemize}
\item پراکندگی زمان و مکان:
\begin{itemize}
\item بین حملات \MShl{حداقل ۱ ماه فاصله} / هر حمله \MShl{حداقل \MSen{24 h}}.
\item درگیری در نواحی مجزایی از سیستم عصبی.
\end{itemize}
\item \MSen{MRI}: پلاک‌های دمیلینه در \MShl{اطراف بطن‌های جانبی}، \MShl{مادهٔ سفید ساب‌کورتیکال}، \MShl{\MSen{Brainstem}}، \MShl{مخچه}، \MShl{نخاع}.
\item مایع \MSen{CSF}: \MShl{افزایش \MSen{IgG}}، ایجاد \MShl{باند الیگوکلونال}.
\begin{itemize}
\item \MShl{\MSen{Pro} و سلول اندکی افزایش یافته}.
\end{itemize}
\item \MShl{پتانسیل برانگیخته}: به ویژه تأخیر برانگیختهٔ بینایی، \MShl{در تشخیص موارد مشکوک}.
\end{itemize}
\section*{سندرم ایزولهٔ بالینی (\MShl{\MSen{CIS}})}
\begin{itemize}
\item بیمارانی که فقط دچار یک حمله گردیده‌اند.
\item در این بیماران تعداد پلاک ارزش تشخیصی و پیش‌بینی‌کننده دارد.
\item \MShl{وجود \MSen{$\geq 3$} پلاک در \MSen{MRI}}: ۷۰٪ طی ۱۰ سال مبتلا به \MSen{MS}.
\end{itemize}
\section*{در چه مواردی در تشخیص \MSen{MS} احتیاط می‌کنیم؟}
\begin{itemize}
\item حملاتی که پراکندگی مکانی وجود نداشته باشد.
\item بیمار هیچ‌وقت علائم بینایی، حسی و اختلال ادراری نداشته باشد.
\item تظاهر با علائم ناشایع: آفازی، کره، اختلال هوشیاری، تشنج.
\item \MSen{MRI} مغز و نخاع طبیعی باشد.
\end{itemize}
\section*{\MSen{MRI} در بیمار \MSen{MS}}
\begin{itemize}
\item \MSen{$T_2$}: پلاک‌ها به شکل مناطق \MShl{\MSen{high Signal}}.
\item \MSen{$T_1$}: \MShl{\MSen{Black hole}}: \MShl{پلاک هایپواینتنس} (\MShl{\MSen{low Signal}}).
\item \MShl{اکثر پلاک‌های مغزی بی‌علامت} ولی \MShl{اکثر پلاک‌های نخاع علامت‌دار هستند}.
\item ارتباطی بین تعداد پلاک و علائم بالینی وجود ندارد.
\begin{itemize}
\item با این حال \MShl{پلاک‌های \MSen{$T_1$} تطابق بالینی بیشتری دارند}.
\end{itemize}
\end{itemize}

\MSpage{۴}{pic04.jpg}
\section*{\MSen{DDX} بیماری \MSen{MS}}
\noindent
ایدز، واسکولیت‌ها (\MSen{SLE}، بهجت، سندرم آنتی‌کاردیولیپین)، بیماری \MSen{Cadasil}، کمبود \MSen{$B_{12}$}، ویروس \MSen{HTLV1}، آنسفالیت حاد دمیلینه‌زان، لکودیستروفی‌ها، سیفلیس \MSen{CNS}، بیماری لایم، سارکوئیدوز، لنفوم مغزی، متاستازها، سندرم \MSen{Susac}، آتاکسی فریدریش، آسیب ایسکمیک پراکنده، تروما.
\section*{پیش‌آگهی \MSen{MS}}
\begin{itemize}
\item خوب: \MSen{♀}، \MShl{شروع بیماری با \MSen{ON}}، شروع قبل از ۴۰ سالگی، کم بودن تعداد حملات در سال‌های اول، \MShl{سیر بهبود ـ عود (\MSen{RRMS})}، میزان ناتوانی اندک در ۵ سال اول.
\item بد: \MSen{♂}، شروع بعد از ۴۰ سالگی، شروع با علائم مخچه‌ای، سیر یکنواخت و پیش‌رونده، شروع با درگیری همزمان چند قسمت از \MSen{CNS}.
\end{itemize}
\section*{درمان \MSen{MS}}
\begin{itemize}
\item مرحلهٔ حاد: \MShl{گلوکوکورتیکوئید}: یک گرم متیل‌پردنیزولون روزانه به مدت ۳–۵ روز \MSen{IV}؛ سپس ۲ هفته پردنیزولون خوراکی.
\item \MShl{پلاسمافرزیس}: در حملات شدید.
\item درمان‌هایی که سیر بیماری را تغییر می‌دهند:
\begin{itemize}
\item اینترفرون بتا: \MShl{کاهش ۴۰٪ عود حملات} / \MShl{برای فرم عودکننده}.
\begin{itemize}
\item عارضه: علائم شبیه به سرماخوردگی.
\end{itemize}
\item گلاتیرامر استات: \MSen{20 mg SC daily}.
\item میتوکسانترون: اگر به داروهای فوق جواب نداد.
\begin{itemize}
\item در موارد پیش‌روندهٔ ثانویه.
\item عوارض قلبی: \MShl{\MSen{HF}}، \MShl{\MSen{CM}}.
\item تزریق هر ۳ ماه یک‌بار: \MShl{قبل از تزریق بررسی قلبی دقیق}.
\end{itemize}
\end{itemize}
\end{itemize}

\MSpage{۵}{pic05.jpg}
\section*{ادامهٔ درمان \MSen{MS}}
\begin{itemize}
\item ناتالیزوماب: اثر آن دو برابر اینترفرون.
\begin{itemize}
\item عارضه: \MShl{لکوآنسفالوپاتی مولتی‌فوکال پیش‌رونده}.
\end{itemize}
\item \MShl{برای \MSen{MS} پیش‌روندهٔ اولیه درمان خاصی وجود ندارد}.
\end{itemize}
\section*{لکوآنسفالوپاتی مولتی‌فوکال پیش‌رونده (\MShl{\MSen{PML}})}
\begin{itemize}
\item مناطق \MShl{مولتی‌فوکال دمیلینهٔ تحت‌حاد} در \MShl{اندازه‌های مختلف} در مغز.
\item عامل بیماری: \MShl{ویروس \MSen{JC}}.
\item \MSen{RF}: \MShl{نقص ایمنی} (\MShl{تقریباً همهٔ بیماران دارند}).
\begin{itemize}
\item ایدز، لوکمی، گیرندهٔ پیوند، بیماری التهابی مزمن (سارکوئیدوز، \MSen{TB}).
\end{itemize}
\item علائم: درگیری لوب اکسیپیتال: \MShl{نقص بینایی} به خصوص همی‌آنوپی هم‌جهت؛ اختلال ذهنی (دمانس)، گیجی، تغییرات شناختی، ضعف شامل همی‌پارزی، پارزی یک اندام، آتاکسی.
\item تشخیص: \MShl{\MSen{MRI} روش ارجح}.
\begin{itemize}
\item \MSen{CSF}: \MShl{طبیعی} / یا \MSen{$pro\uparrow$}، \MSen{$IgG\uparrow$} / در \MSen{$\tfrac{1}{4}$} موارد پلئوسیتوز.
\item \MSen{DNA}: کشف \MSen{DNA} ویروس در \MSen{CSF}.
\end{itemize}
\item درمان: \MShl{درمان مؤثری وجود ندارد}.
\begin{itemize}
\item \MShl{سیتارابین} و \MShl{سیدوفوویر}.
\end{itemize}
\end{itemize}
\section*{درمان‌های علامتی \MSen{MS}}
\begin{itemize}
\item اسپاستیته: باکلوفن، تیزانیدین، دانترولن.
\item خستگی: آمانتادین، متیل‌فنیدیت، مودافینیل.
\item اختلال نعوظ: سیلدنافیل.
\item درد: کاربامازپین، گاباپنتین، نورتریپتیلین.
\item افسردگی: فلوکستین، آمی‌تریپتیلین.
\end{itemize}
\endgroup
\ifMSStandalone
  \expandafter\enddocument
\fi
